\chapter{Configuration}
\label{ch:config}

Secrets live only in a \code{.env} file (permissions \code{0600}, never committed, never
logged). Non-secret settings live in a TOML config file the UI can re-read at runtime.

\section{smart-me authentication}
\label{sec:smartme-auth}

smart-me offers several authentication mechanisms. \prog uses the one designed for
headless devices.

\subsection{The available mechanisms}
\begin{itemize}
  \item \textbf{OAuth “Device (Client Credentials)” app} — described by smart-me as
        \emph{“a replacement for basic authentication in embedded devices\dots\ supports
        only the OAuth flow client credentials.”} This is a non-interactive
        machine-to-machine flow — exactly what a 24/7 daemon needs.
  \item \textbf{OAuth Confidential / Public apps} — use interactive flows
        (authorization-code~+~PKCE, device-code, implicit) that require a human to log in
        and consent. Not suitable for an unattended bridge.
  \item \textbf{HTTP Basic auth} — the smart-me account username and password.
\end{itemize}

\subsection{What \prog uses}
\begin{description}[style=nextline]
  \item[Primary — OAuth2 Client Credentials] The bridge exchanges a Client~ID and Client
        Secret for a bearer token, then calls the API with
        \code{Authorization: Bearer <token>}, refreshing the token on expiry.
  \item[Fallback — HTTP Basic] Account user/password. \textbf{Designed, implemented in the
        client library, and not reachable from the environment} --- see the warning below.
\end{description}

\begin{warning}
\textbf{The Basic fallback cannot currently be selected.} \code{smart-me-client} implements
it, but the binary builds OAuth client-credentials unconditionally and never constructs
Basic credentials, so \code{SMARTME\_BASIC\_USER} and \code{SMARTME\_BASIC\_PASSWORD} have
no effect whatever. Worse for anyone relying on the word \emph{fallback}: the OAuth pair is
\emph{required}, so setting the Basic pair \emph{instead} of it makes \prog refuse to start
with \code{missing environment variable SMARTME\_CLIENT\_ID}. Treat client credentials as
the only way in until this is either wired up or dropped from the design.
\end{warning}

\begin{note}
This decision — client credentials primary, Basic fallback, TLS mandatory — is recorded
in \code{docs/adr/0009-smartme-auth-client-credentials.md}. It supersedes the earlier
\code{Authorization: ApiKey} assumption.
\end{note}

\subsection{Configuring it}
Create a \emph{Device (Client Credentials)} OAuth app in the smart-me portal
(\href{https://api.smart-me.com/swagger}{api.smart-me.com}), copy its Client~ID and
secret, and set them in \code{.env}:

\begin{lstlisting}[language=bash]
# smart-me credentials -- LOCAL ONLY, never commit (.env is gitignored, perms 0600)
SMARTME_CLIENT_ID=your-client-id
SMARTME_CLIENT_SECRET=your-client-secret

# HTTP Basic fallback -- RESERVED, and NOT read by the bridge today.
# Setting these has no effect; the OAuth pair above is still required.
# SMARTME_BASIC_USER=your-account-user
# SMARTME_BASIC_PASSWORD=your-account-password
\end{lstlisting}

\begin{warning}
If the client secret is lost or leaked, smart-me does not allow rotating it in place: the
OAuth app must be deleted and recreated. All smart-me traffic is over TLS and hard-fails
if TLS is unavailable.
\end{warning}

\subsection{Scopes: what an app is allowed to do}
\label{sec:smartme-scopes}

An app's power comes from the \emph{claims} granted to it when it is created, not from the
\code{scope} parameter sent at token exchange. A token request may only ask for claims the
app already holds; asking for more is refused.

\begin{warning}
Do not take the scope list from the OpenID discovery document
(\code{/.well-known/openid-configuration}). It advertises only \code{openid} and
\code{offline\_access}, which is incomplete to the point of being misleading — the real
vocabulary is twelve claims. Read it instead from the OAuth2 security scheme in
\code{https://api.smart-me.com/swagger/v1/swagger.json}, or live from
\code{GET /oauth/claimtree}, which also returns the hierarchy.
\end{warning}

The claims form a tree, and \textbf{a parent implies its children}. Granting
\code{device.readswitch} therefore also grants \code{device.read} — a switching app is
always also a reading app, and there is no claim that permits switching \emph{without}
reading.

\begin{center}
\begin{tabular}{@{}>{\ttfamily}l p{7.6cm}@{}}
\toprule
\normalfont\textbf{Claim} & \textbf{Grants} \\
\midrule
device.readwrite      & Read, control outputs, and add/update values \\
\quad device.readswitch & Read and control the outputs (the relay) \\
\quad\quad device.read  & Read device data only \\
user.readwrite        & Read and write the user settings \\
\quad user.read       & Read the user settings \\
auditlog.read         & Read the audit log \\
realtimeapi.readwrite & Read and write realtime-API registrations \\
\quad realtimeapi.read & Read realtime-API registrations \\
partner.write         & Read and write a partner's devices \\
\quad partner.read    & Read a partner's devices \\
partner.updatefirmware & Update the firmware of all a partner's devices \\
oauthapp.write        & Create and change OAuth applications \\
\bottomrule
\end{tabular}
\end{center}

\subsection{Creating a read-only app (what \prog needs)}

\prog only ever reads. Grant it \code{device.read} and nothing else: a bridge that cannot
switch a relay cannot switch one by accident, and a leaked secret cannot be used to act on
the installation.

\begin{enumerate}
  \item In the smart-me portal, create an OAuth application of type
        \emph{Device (Client Credentials)}.
  \item Give it a name that says where it runs (for example \code{gateway}).
  \item Grant the single claim \code{device.read}.
  \item Copy the Client~ID and secret into \code{.env} as shown above.
\end{enumerate}

\paragraph{Verifying that it really is read-only.} This can be proved without attempting a
write, because the token endpoint distinguishes the two failure modes:

\begin{lstlisting}[language=bash]
# Ask for a claim the app should NOT have.
curl -s -X POST https://api.smart-me.com/oauth/token \
  -d grant_type=client_credentials \
  -d client_id="$ID" -d client_secret="$SECRET" \
  -d scope=device.readswitch
\end{lstlisting}

\begin{description}[style=nextline]
  \item[\code{invalid\_request} — \emph{“This client application is not allowed to use the
        specified scope.”}] The claim exists and the app does not hold it. This is the
        expected answer for a read-only app, and it is a positive result.
  \item[\code{invalid\_scope}] The claim name does not exist. You have made a typo, and the
        test has proved nothing.
  \item[A token is returned] The app \emph{can} switch outputs. It is not read-only.
\end{description}

\begin{note}
Run this against the credentials actually deployed, not against the ones you meant to
deploy. The distinction between the two error codes is what makes the check meaningful: a
test that cannot tell “refused” from “misspelled” always passes.
\end{note}

\subsection{Creating an app that can command a meter relay}

\begin{warning}
\prog does not use this, and does not implement any command path. The bridge publishes
measurements and answers Sparkplug \code{NCMD} \code{Node Control/Rebirth}; it never writes
to smart-me. This section documents the smart-me side so that the capability can be
enabled deliberately — it is not a description of current behaviour. The corresponding
Sparkplug \code{DCMD} question is tracked as issue~\#38.
\end{warning}

Create a second, \emph{separate} Device app granted \code{device.readswitch}. Keep it
separate from the bridge's read-only app: the two have different blast radii, and only one
of them is worth protecting carefully.

\paragraph{1 — Find the device's actions.} Not every meter has a switchable output. Ask
before assuming:

\begin{lstlisting}[language=bash]
curl -s -H "Authorization: Bearer $TOKEN" \
  https://api.smart-me.com/actions/<device-uuid>
\end{lstlisting}

A meter with no relay returns \code{[]}. A meter with one returns its actions, and the
\code{ObisCode} is the identifier you will command:

\begin{lstlisting}
[{"Name":"Relais","ObisCode":"63000C0102FF","ActionType":0}]
\end{lstlisting}

\paragraph{2 — Send the command.}

\begin{lstlisting}[language=bash]
curl -s -X POST https://api.smart-me.com/actions \
  -H "Authorization: Bearer $TOKEN" \
  -H 'Content-Type: application/json' \
  -d '{"deviceID":"<device-uuid>",
       "actions":[{"obisCode":"63000C0102FF","value":1}]}'
\end{lstlisting}

\begin{note}
\code{PUT /Devices/\{id\}} also switches an output and still works, but smart-me's own
specification says of it: \emph{“For new implementations please use the \code{actions}
command.”} Prefer \code{POST /actions}.
\end{note}

\begin{tip}
\code{GET /actions/\{id\}} needs only \code{device.read}, so you can inventory which meters
have a relay using the bridge's existing read-only credentials, before creating any app
that can actually operate one.
\end{tip}

\section{MQTT broker connection}
\label{sec:broker-config}

\prog connects to an \emph{external} broker; it does not embed one. Host and port are
deployment-specific, so they belong in \code{.env} alongside the credentials and never in
a committed file:

\begin{lstlisting}[language=bash]
# MQTT broker -- LOCAL ONLY, never commit
SMARTME_BROKER_HOST=mqtt.example.lan   # host or IP of your broker
SMARTME_BROKER_PORT=1883               # optional; defaults to 1883
\end{lstlisting}

A port that is not a number makes the bridge refuse to start. It does not fall back to
\code{1883} silently --- connecting to a different broker than intended is precisely the
kind of quiet wrong answer this project avoids.

\subsection{Authentication and transport security}

Both are \textbf{planned and optional} --- the broker connection may be secured with TLS
and/or authentication, or left plain, per configuration (FR43, NFR16). Neither is
implemented in the current build.

\begin{warning}
As of this build the broker connection is \textbf{plain and unauthenticated}. \prog
therefore expects its broker to sit on a \textbf{trusted internal network}, reachable only
from hosts you control. Traffic to the \emph{smart-me cloud} is unaffected and remains
TLS-only, hard-failing if TLS is unavailable --- the two paths use different TLS stacks
and only the broker one is disabled.
\end{warning}

\subsubsection*{Why optional rather than mandatory}

A homelab broker is usually shared. It commonly serves devices --- older sensors,
closed-firmware appliances --- that cannot present credentials at all, so a bridge that
\emph{required} authentication would either force the broker into a configuration its
other clients cannot meet, or force a second broker into existence. Enabling
authentication is a decision about the whole broker, not about this bridge, and \prog
follows it rather than driving it.

\subsubsection*{Reserved settings}

The variable names are reserved now so that enabling authentication later is a
configuration change rather than a migration:

\begin{lstlisting}[language=bash]
# Broker authentication -- reserved, not yet honoured by the bridge
# SMARTME_BROKER_USER=
# SMARTME_BROKER_PASSWORD=
\end{lstlisting}

\begin{note}
Broker credentials will follow the same discipline as the smart-me ones: \code{.env} only,
never logged, never in a committed file (NFR16).
\end{note}

\subsubsection*{Authentication and TLS are worth landing together}

MQTT carries username and password \emph{in the clear} inside the CONNECT packet. On a
plain connection, enabling authentication protects against an accidental or careless
client on the same network, but not against anyone able to observe the traffic --- they
would simply read the credentials as they go past. On a trusted LAN that is still a real
improvement; treating it as confidentiality is not.

Broker TLS is currently blocked upstream: the MQTT client's default TLS stack pinned a
transitive dependency carrying unfixed advisories that could not be upgraded past its own
requirement, so it was disabled rather than shipped unused. When that upgrade lands, TLS
and authentication are best enabled in the same change.

\subsection{Sparkplug identity}

The group and edge-node identifiers are \textbf{required}. There is no default, and
\prog refuses to start without them:

\begin{lstlisting}[language=bash]
SMARTME_GROUP_ID=<your group>
SMARTME_NODE_ID=<your node>
\end{lstlisting}

\begin{warning}
These form the topic grammar, which is a \emph{versioned contract} (chapter~\ref{ch:contract}).
Changing them after a SCADA has begun historising the tags breaks the continuity of its
stored curves --- the history stays under the old topic and the new one starts empty.
Choose them before the first production run, not after.

\textbf{And a Sparkplug host persists what it discovers.} The group you name becomes a
folder in the host's tag tree that outlives the process and has to be deleted by hand ---
and deleting those tags also discards their alarm and history configuration. Publishing
into the wrong namespace is not undone by restarting with better settings, which is why
these have no default at all: they defaulted to \code{Site} and \code{Bridge} until
2026-07-31, so a \prog started with no configuration published straight into the namespace
most likely to be the production one.
\end{warning}

\section{Meter-to-topic mapping}
\stub{default scheme, editing, first-run mapping confirmation.}

\section{Poll interval, logging, and retention}

\stub{configurable poll interval.}

\subsection{Logging to a file}

The console always receives every line. File logging is \emph{additional} and opt-in: with
\code{SMARTME\_LOG\_DIR} unset there is no log file at all, and \prog behaves exactly as it
did before file logging existed.

\begin{lstlisting}[language=bash]
# Turns file logging ON. Written as smartme_mqtt.<YYYY-MM-DD>.log, rotated daily.
SMARTME_LOG_DIR=/data/logs

# How many rotated files to keep. Default 7 (one week).
SMARTME_LOG_KEEP=7
\end{lstlisting}

\code{SMARTME\_LOG\_KEEP} must be a positive whole number. A value that is not one is
refused rather than guessed at: \prog says so on stderr and keeps 7.

\begin{note}
File logging can never stop the bridge. If the directory cannot be created or the file
cannot be opened, logging degrades to console-only and prints the reason on stderr. In a
container the usual cause is the bind-mount's ownership --- the directory must belong to the
uid the container runs as (10002); see chapter~\ref{ch:install}.
\end{note}

\begin{warning}
Under pressure the \emph{file} loses lines and the console does not. A saturated write
buffer drops log lines rather than blocking the bridge, because a slow disk must never
stall publishing. So \code{docker compose logs} stays complete when the file does not. The
number of lost lines is reported when \prog shuts down, so the loss is never silent --- if
you are reconciling a file against the console, that line is the reason they differ.
\end{warning}

\subsection{Log level}

\code{RUST\_LOG} is the standard \code{tracing} filter and applies to the console and the
file alike. It defaults to \code{info}, deliberately: the library's own default is
\code{error}, which would hide the ignored-command traces, the QoS-downgrade warning and
every traced drop --- an operator would read an empty log as ``nothing happened''.

\begin{lstlisting}[language=bash]
# Default.
RUST_LOG=info

# Raise one subsystem without drowning in the rest.
RUST_LOG=info,smartme_bridge::adapters::mqtt_driver=debug
\end{lstlisting}
